\chapter{从图像生成到多模态世界}
\label{ch:applications}

\readingguide{前两章告诉我们怎样学习一个把噪声变成数据的场，第三章告诉我们怎样训练和积分这个场。本章要问：这个场由什么网络表示，又怎样接收人的意图？我们先拆开图像生成系统，再沿着 U-Net、隐空间、DiT、视频和统一多模态模型的路线前进。阅读时始终区分概率路径、网络结构、条件接口与数据，它们共同决定系统能力，却不是同一件事。文献与产品信息核查截至 2026 年 10 月 9 日；2026 年的研究例子明确标注为预印本，产品公开演示不作为未披露架构的证据。}

\section{先拆开一套生成系统}
\label{sec:application-system}

看到一张清晰的生成图，最容易产生的印象是：模型内部一定藏着一个直接画图的巨大函数。实际系统通常更像一条分工明确的生产线。文字先被编码成条件；随机变量提供尚未确定的细节；一个神经网络反复预测噪声、干净样本、score 或速度；采样器把预测变成状态更新；最后，解码器可能把隐变量变成像素。用户看到的是最终图像，训练者需要同时照顾这条链上每一个环节。

可以把一套系统写成一个配置组
\[
\mathcal S=(\text{表示空间},\text{概率路径},\text{预测目标},
\text{网络},\text{条件},\text{采样器},\text{数据}).
\]
这个写法有一个实际好处：它阻止我们把性能变化全部归因于论文标题里最醒目的那个名词。把 U-Net 换成 Transformer，改变的是函数的参数化；把噪声预测换成直线路径速度回归，改变的是训练目标和概率路径；把像素换成 VAE 隐变量，改变的是网络处理的随机变量；把短标题换成详细描述，改变的是监督信号。它们能够组合，但没有哪一项独自保证生成质量。

图像生成网络的输入和输出往往形状相同。例如，输入是一个含噪的 $h\times w\times C$ 张量，输出也是同样大小的噪声或速度张量。这个外观容易让人以为生成就是图像到图像的回归。关键区别在于，输入的噪声水平也作为条件，而且同一个网络要覆盖整条概率路径。它学到的是许多分布之间的局部迁移规律，采样器才把这些局部规律连接成一张完整图。可以回看第\ref{ch:math}章的式\eqref{eq:reverse-sde}、\eqref{eq:pfode}和\eqref{eq:cfm}，把其中抽象的场理解成下文讨论的网络输出。

\begin{intuitionbox}{网络负责方向，采样器负责行走}
U-Net 和 DiT 回答的是“怎样用有限参数表示这个高维场”；DDPM、DDIM 和概率流 ODE 回答的是“沿着什么随机或确定性规则更新状态”。一张图是否清晰，既取决于方向是否准确，也取决于步长、条件和表示空间是否合适。
\end{intuitionbox}

这也解释了为什么软件中的一个模型名称，未必对应一种固定数学定义。一个产品系列可以在不同代际中更换主干网络、训练路径和文本编码器。阅读论文时，应当落实到具体版本、具体训练设置。我们下面使用历史工作作为概念的锚点，不把某个版本的实现推广为整个品牌永远不变的性质。

\section{U-Net 为什么适合学习图像中的场}
\label{sec:application-unet}

\subsection{从局部纹理到整体布局}

一张图像同时包含多个尺度的信息。头发的细纹只占几个像素，人脸的比例需要几十个像素，人物和背景的布局则跨越整张画面。如果每次去噪只看很小的邻域，模型可能把局部画得精致，却无法协调左右手、物体数量和远处背景。反过来，如果网络一开始就把画面压缩得很粗，细节的位置又可能丢失。

U-Net 以一条下降路径和一条上升路径回应这个矛盾。下降路径逐渐减小空间分辨率、增加通道数，把局部证据汇总成较大的语义结构；上升路径逐渐恢复空间分辨率，把整体决定落实到局部。对应尺度之间的跳跃连接直接传递浅层特征，因此深层网络不必凭空重建所有精细位置信息。这里的“U”描述拓扑形状，既不是某一种噪声日程，也不指定采样随机性。

图\ref{fig:application-unet}给出的是原创教学示意，不是某篇论文的逐层复刻。含噪输入经过不同尺度，最终得到一个同形状的预测。中间的跳跃连接传递的是特征，而不是自动保留了正确答案。尤其在高噪声阶段，浅层特征本身充满噪声，模型仍需根据时间条件决定如何使用它们。

\begin{figure}[htbp]
\centering
\begin{tikzpicture}[x=1cm,y=1cm, every node/.style={font=\small},
  block/.style={draw,rounded corners,align=center,minimum width=2.5cm,minimum height=0.8cm},
  arr/.style={->,thick}]
\node[block] (a) at (0,0) {高分辨率特征\\含噪输入};
\node[block] (b) at (0,-1.5) {中分辨率特征};
\node[block] (c) at (0,-3) {低分辨率特征};
\node[block] (d) at (6,-3) {低分辨率融合};
\node[block] (e) at (6,-1.5) {中分辨率融合};
\node[block] (f) at (6,0) {高分辨率预测\\噪声或速度};
\draw[arr] (a)--node[left]{下采样}(b);
\draw[arr] (b)--(c);
\draw[arr] (c)--node[above]{大范围交互}(d);
\draw[arr] (d)--(e);
\draw[arr] (e)--node[right]{上采样}(f);
\draw[arr,dashed] (a)--node[above]{跳跃连接}(f);
\draw[arr,dashed] (b)--node[above]{跳跃连接}(e);
\node[align=center] at (3,-4.1) {时间嵌入告诉每一层当前噪声水平；\\文字或类别条件告诉每一层要生成什么。};
\end{tikzpicture}
\caption{U-Net 的尺度分工。虚线传递对应尺度的特征；图中省略了残差块和具体注意力位置。}
\label{fig:application-unet}
\end{figure}

卷积提供了一个很有用的归纳偏置：同一种局部规则可以在画面的不同位置重复使用。一个边缘检测器不必为左上角和右下角分别学习。它特别适合纹理与局部几何，但纯卷积的长距离交互要经过多层传播。于是扩散 U-Net 通常还引入注意力，让相距较远的特征直接交换信息。注意力不是专门为语言设计的，它是一种按内容建立连接的方法，术语解释见\hyperref[gl:attention]{注意力条目}。

\subsection{时间如何进入网络}

把一幅几乎干净的图和一幅接近纯噪声的图交给同一个网络，却不说明噪声水平，会使回归目标变得含混。网络看到一个弱边缘时，不知道它是应该保留的图像结构，还是应当消除的噪声偶然性。因此，时间 $t$ 或离散步数 $k$ 通常先变成一个向量，再被送到许多残差块。

一个常见的构造是先用正弦、余弦形成不同频率的编码，再经过小型全连接网络。这样的编码把一个标量变成多个平滑变化的坐标，有助于网络区分相近的噪声水平。时间信息可以直接加到特征上，也可以产生特征归一化后的缩放和偏移。若 $h$ 是某层特征，后一种方式可以示意为
\[
\widetilde h=\alpha(t,c)\odot\operatorname{Norm}(h)+\beta(t,c).
\]
这里 $\odot$ 是逐元素乘法，$\alpha$、$\beta$ 由条件网络产生。不要把本式的 $\alpha(t,c)$ 与 DDPM 的 $\alpha_k=1-\beta_k$ 混为一谈：前者是神经网络的特征调制参数，后者是概率路径的系数。它们只是在不同文献中恰好使用了同一个希腊字母。

可以把特征调制理解为网络内部的一组工作模式。同样一层卷积，在大噪声时侧重粗略结构，在小噪声时侧重边缘和材质。这个比喻并不意味着每个通道都能得到清晰的人类语义解释；它说明的是条件可以在许多层持续改变计算，而不只是作为开头的一个标签出现。

\subsection{类别条件走向文字条件}

类别条件只需说明“这张图属于哪一类”。文字条件却包含属性、关系、数量和动作。一条提示词“左侧是一只红杯子，右侧是一只蓝杯子”，如果只压缩成一个整体向量，红、蓝与左右的绑定就可能不稳定。保留文字的 token 序列，允许不同图像位置查询不同文字片段，是更细致的接口。

令图像特征产生查询 $Q=H W_Q$，文字特征产生键和值 $K=CW_K$、$V=CW_V$。交叉注意力写为
\begin{equation}
\operatorname{CrossAttn}(H,C)
=\operatorname{softmax}\!\left(\frac{QK^\top}{\sqrt{d_k}}\right)V.
\label{eq:application-crossattention}
\end{equation}
矩阵的一行对应某个图像位置或图像块。该行通过文字键的相似度，对文字值进行加权汇总。$d_k$ 是单个注意力头中键的维数，除以其平方根是为了避免内积幅度随维数增长过快。softmax 是按行计算的归一化指数函数，因此每行权重和为一。

这个机制并不保证“红色”一定绑定到正确的杯子。训练只通过最终回归误差监督注意力，数据描述也可能只说“桌上的杯子”，没有提供所有关系。注意力提供了表示关系的通道，关系是否被学会还取决于数据、编码器和优化。这个区别非常重要：一种结构能表达某个能力，与训练后可靠拥有这个能力，是两种不同结论。

\section{文字到图像：条件比人们想象的更重要}
\label{sec:application-textimage}

\subsection{ADM：先建立一个高质量图像先验}

Dhariwal 与 Nichol 的 ADM 工作通过 U-Net 结构实验和分类器引导，说明扩散模型在当时的图像基准上可以与强 GAN 系统竞争。论文的关键教学价值，是把高质量图像先验与引导机制分开研究，而不是把优势归结为“多做几步去噪”。\cite{a_adm}

用第\ref{sec:guidance}节的语言说，类别引导在无条件 score 上增加 $\grad_x\log p(c\mid x_t)$。梯度告诉我们：在当前噪声水平附近，朝哪个方向移动更容易被识别成所需类别。引导系数增大时，样本可能更典型，但也更集中。图像“更像一只狗”不等于狗的全部姿态、品种和背景都被更好地覆盖。第\ref{ch:concepts}章中学习分布的目标，到这里变成了一个很现实的质量与覆盖率权衡。

这里也可看出分类器的额外要求。它不能只识别干净图像，还必须在各种噪声水平下给出有意义的梯度。如果分类器在高噪声处不稳定，错误方向会进入整个采样过程。这推动了不另训噪声分类器的条件方法发展。

\subsection{GLIDE：让文本参与每次去噪}

GLIDE 把文字条件引入扩散图像生成，比较了 CLIP 引导与\term{无分类器引导}{gl:cfg}，并展示文字引导的图像补全。其研究结果支持了把有条件与无条件预测组合起来的实用路线；公开小模型与论文大模型的规模和训练条件应分别阅读。\cite{a_glide}

无分类器引导之所以容易移植，源于它只要求网络在训练时有时看见条件、有时看不见条件。推理时，两次预测的差给出一个与条件相关的方向。这个做法可以对噪声、score 或恰当参数化的速度进行线性组合，但必须保持训练路径和参数换算一致，具体公式见式\eqref{eq:cfg}。它把“更听话”变成可调节的采样操作，并不自动补足训练数据中缺失的概念。

如果模型从未学会某种空间关系，把引导强度继续增大，可能只是强化一种错误解释。实践中的饱和、过强对比度和形态扭曲，提醒我们：引导是改变已有场，而不是向网络注入全新的知识。

\subsection{Imagen：语言理解也是图像生成的瓶颈}

Imagen 使用预训练大型文本编码器和级联扩散模型，并通过实验强调文本表征对画面质量与图文一致性的作用。论文还引入包含复杂提示的 DrawBench，帮助评估数量、属性和组合关系。\cite{a_imagen}

一个冻结的文本编码器为什么仍然重要？因为去噪网络并不直接看到自然语言字符串，它只能看到编码器输出的向量。若两条本应不同的提示在表征中几乎重合，后面的生成器再强，也难以稳定区分。反过来，语言编码器保留了语义差异，图像网络才有机会把差异落实到空间布局。

级联生成体现了另一种分工。基础模型先生成低分辨率图像；超分辨率模型在其基础上生成更细的图像。数学上，这可以示意为
\[
p(x^{\mathrm{high}}\mid c)
=\int p(x^{\mathrm{high}}\mid x^{\mathrm{low}},c)
        p(x^{\mathrm{low}}\mid c)\,\dd x^{\mathrm{low}}.
\]
低分辨率结果已经承担了构图的很大部分，因此高分辨率模型主要解决在给定结构下补足细节的问题。它能减轻一次生成整张大图的困难，却也可能继承上一级的计数或布局错误。把三只杯子放大，通常不会自动变成正确的四只杯子。

\subsection{DALL-E 2：语义隐变量也是一种生成对象}

DALL-E 2 的论文把图像生成分成文字到 CLIP 图像表征的先验，以及从该表征到图像的解码器。先验研究了自回归和扩散方案，解码阶段使用扩散生成。它展示了一种与空间 VAE 隐变量不同的语义中间层。\cite{a_dalle2}

CLIP 指对比式语言图像预训练：文字与图像分别经过编码器，配对样本的表征被拉近，不配对样本的表征被区分。它因此提供一个适合比较图文语义的空间。\cite{a_clip}CLIP 图像向量不是一张缩小后的图片，也不是像素的无损编码。

为了理解这一区别，考虑两个问题：“画面总体表达什么？”与“每一个像素在哪里？”语义表征更适合前一个问题，保留整体概念和图文对齐；空间隐变量更接近后一个问题，保留可解码的局部结构。两者都叫 latent，却服务于不同接口。把“用了隐空间”直接理解成“用了 Stable Diffusion 那样的压缩图像”，会错过模型的真正设计。

中间层越抽象，语义控制可能越便利，但对精确文字、细小位置和对象数量的保留也可能越困难。这个取舍会反复出现：生成系统需要既懂意义又记得细节，而一种紧凑表征很难免费满足二者。

\section{隐空间扩散：把算力花在更合适的变量上}
\label{sec:application-latent}

\subsection{从像素分布到压缩表征分布}

高分辨率像素数量很大，其中相当一部分变化来自局部纹理、采样噪声和细微颜色。一个生成器若要在每个采样步都处理所有像素，计算量会很快增加。隐空间扩散模型（Latent Diffusion Model，LDM）的重要选择，是先学习一个图像自编码器，再在它的空间隐变量上学习扩散先验，并用交叉注意力接入条件。\cite{c_rombach2022ldm}

记编码器为 $E_\phi$，解码器为 $D_\psi$，压缩后的干净隐变量为 $\ell_0$。实际编码器可以输出后验分布的参数，而不是唯一的确定向量；为了突出系统结构，暂时用 $\ell_0=E_\phi(x_0)$ 记一次编码结果。隐空间训练对 $\ell_0$ 加噪：
\[
\ell_t=a(t)\ell_0+b(t)\eps.
\]
第二章的去噪推导仍然适用，只是随机变量已经从像素换成隐变量。生成时先从噪声得到 $\widehat\ell_0$，再得到图像 $D_\psi(\widehat\ell_0)$。在随机解码器的写法下，总分布为
\begin{equation}
p_{\theta,\psi}(x\mid c)
=\int p_\psi(x\mid\ell)\,p_\theta(\ell\mid c)\,\dd\ell.
\label{eq:application-latent-model}
\end{equation}
若解码器是确定性的，就把 $p_\psi(x\mid\ell)$ 理解为集中在 $D_\psi(\ell)$ 处的点质量。该积分描述的是解码器对隐变量分布的映射，术语可对照\hyperref[gl:latent]{隐空间条目}和\hyperref[gl:vae]{VAE 条目}。

注意，VAE 中的先验正则不意味着真实编码结果严格服从标准高斯。图像经编码后形成的聚合分布仍然复杂；扩散模型正是来拟合这个分布。否则，若标准高斯直接经解码器已经能生成足够好的图像，就没有再训练昂贵隐空间先验的必要。

\subsection{为什么压缩并非越强越好}

设原图是 $H\times W\times 3$，编码器把空间边长缩小 $f$ 倍，隐通道数为 $C_\ell$。元素数之比为
\[
\frac{(H/f)(W/f)C_\ell}{3HW}
=\frac{C_\ell}{3f^2}.
\]
这个式子只是存储元素数的比较，并不直接等于训练速度提升：通道数、网络宽度、注意力和解码成本都参与实际计算。但它清楚显示了空间压缩为何有吸引力。例如一个教学设定 $f=8$、$C_\ell=4$，隐变量元素数是像素元素数的 $1/48$。这里是计算示例，不是对所有模型配置的断言。

代价是压缩会丢信息。若一串小字号文字在编码和重建之后已经模糊，隐空间生成器再精确，也不能从这个接口稳定恢复每个字形。它可能学到一个看似合理的纹理补偿，却没有受监督保留原来的文字。这是表示瓶颈，而不是采样器做得不够好。

\begin{examplebox}{先判断错误来自哪里}
假设一个模型总把商品包装上的小字画错。先把真实包装图送进编码器再解码。如果重建图已经错了，优先检查压缩率和自编码器；如果重建正确而生成错误，再检查文字条件、数据和生成先验。把所有错误都交给“更多去噪步数”处理，常常找错了层级。
\end{examplebox}

自编码器的目标也未必只是像素均方误差。它可以结合重建、\term{感知损失}{gl:perceptual}和适当的正则，使重建图保留人们在视觉上在意的结构。感知损失通常在某个特征空间比较图像，而不是逐像素比较。它有助于保持视觉锐度，但优化的“相似”取决于该特征空间，因此仍需按具体任务验证。

\subsection{Stable Diffusion 是一套可组合的接口}

以公开的 Stable Diffusion v1 系统为例，它把文本编码器、隐空间去噪 U-Net、自编码器和采样器连接起来。\cite{a_sdcard}理解这个组合比记住某一代模型的层数更有长期价值。文字编码器告诉模型意图，去噪主干学习条件隐变量分布，自编码器负责像素接口，采样器负责数值更新。只要接口保持相容，其中某个组件就可以改进而不必重训所有组件。

这解释了生态里为什么会出现许多不同采样器、条件插件与个性化微调。但“可连接”不等于任意组合都有效。更换文本编码器会改变特征维数和语义分布；更换 VAE 会改变隐变量尺度和几何；更换预测目标会改变采样器读取输出的方式。兼容性既是形状问题，也是训练分布问题。

SDXL 在隐空间路线中进一步研究更大的去噪网络、文本条件和分辨率相关条件，并包含基础模型与细化模型的安排。它特别适合作为一个提醒：高分辨率表现来自网络、数据、训练任务和条件的联合工程，而不是单独增加画布尺寸。\cite{a_sdxl}

我们可以亲自理解分辨率条件为什么有用。同一张照片，裁成正方形和保留原宽高比，画面中对象的位置会不同。如果训练时只告诉网络文字，不告诉它原尺寸和裁剪位置，这些不同结果就被混在同一个条件分布里。尺寸、宽高比和裁剪信息作为附加条件，可以让模型解释一部分原先看似随机的布局差异。

\section{DiT：把场学习变成 token 之间的协作}
\label{sec:application-dit}

\subsection{图像块为什么能成为 token}

Transformer 不要求输入必须是文字。它需要的是一个向量序列，以及识别这些向量位置或关系的办法。把隐变量切成互不重叠的小块，将每块展开并线性投影，就得到一串\term{token}{gl:token}。这里 token 是计算中的一个序列元素，它可以对应词片段，也可以对应\term{图像块}{gl:patch}，不能一律理解为文字词语。

若隐变量为 $h\times w\times C_\ell$，块边长为 $p$，并且 $h,w$ 可被 $p$ 整除，则图像 token 数量为
\begin{equation}
N=\frac{hw}{p^2}.
\label{eq:application-token-count}
\end{equation}
每个块经投影变成宽度 $D$ 的向量。位置编码告诉网络它来自哪里，否则一个向量序列在置换之后，网络难以知道“左边”和“右边”。Transformer 块让这些 token 通过自注意力交换信息，然后通过逐 token 的前馈网络做特征变换。最后，预测头把向量还原成每个块中的噪声或速度值，并重新拼成张量。

原始 DiT（Diffusion Transformer，扩散 Transformer）论文在隐空间中用 Transformer 替代扩散 U-Net，系统比较条件注入方法和模型计算规模。其典型设计使用自适应层归一化及零初始化残差调制；主要实验是类别条件生成，不应直接说原始 DiT 就是一套完整文字到图像系统。\cite{a_dit}

\subsection{条件调制与稳定优化}

Transformer 的残差更新可以示意为
\begin{align}
\widetilde H&=(1+\alpha_c)\odot\operatorname{LN}(H)+\beta_c,\\
H'&=H+\gamma_c\odot\operatorname{Attn}(\widetilde H).
\label{eq:application-adaln}
\end{align}
这里 $\operatorname{LN}$ 是层归一化，条件网络从时间和条件产生 $\alpha_c,\beta_c,\gamma_c$，公式省略了多头拆分和后续前馈子层。缩放、偏移决定注意力看到的特征，门 $\gamma_c$ 决定残差更新的强度。零初始化某些调制输出，可让训练开始时的残差分支接近零，深层网络先接近容易优化的恒等映射，再逐步学习复杂更新。

这个设计连接了数学与优化：理论上我们需要一个准确的时间相关向量场，实际上还需要训练从随机参数出发能稳定到达它。损失函数正确不代表优化会顺利，网络中的归一化和初始化属于这条桥梁的一部分。

\subsection{为什么“大一点”可能更好，也可能更贵}

一层标准全注意力中，注意力矩阵有 $N^2$ 个位置对，主要相关运算规模约为 $N^2D$；线性投影和前馈网络的主要规模约为 $ND^2$。忽略常数后，一次网络调用的计算可以写成
\[
\operatorname{Cost}\approx L\bigl(c_1N^2D+c_2ND^2\bigr),
\]
其中 $L$ 是层数，$c_1,c_2$ 是由具体结构决定的常数。实际高效实现可以减少存储开销、利用硬件并行，但不会凭空消除所有 token 交互成本。这个估算用来理解尺度，不用于预测某张显卡的精确运行时间。

减小块大小 $p$ 会增加 $N$。若从 $p=4$ 改成 $p=2$，token 数变为四倍，$N^2$ 部分变为十六倍。与此同时，空间描述更精细，模型可以在更细的粒度上交换信息。增加深度、宽度和 token 数，并非三种完全等价的扩容：它们提供不同的计算能力，也改变不同的瓶颈。

\begin{keybox}{规模规律是一种受条件约束的经验规律}
同一训练体系中，更多计算可能带来更低验证误差和更好样本；这不能推出任意数据、任意网络、任意预算下，参数更多都更好。比较时至少应报告数据、训练计算、采样次数和评价设置。
\end{keybox}

U-Net 的多尺度结构和卷积归纳偏置使局部视觉规律容易学习；DiT 的统一序列接口使长距离交互和多模态扩展自然。两者的区别不该简化成“旧网络”和“新网络”。在预算有限、任务较窄或局部结构特别重要时，卷积式设计仍有合理性；在大规模预训练、复杂条件和统一接口中，Transformer 展现了另一组优势。

\subsection{MM-DiT：文字不只是一份静态查询资料}

SD3 对应的研究结合直线路径的流训练、噪声水平采样研究和多模态 Transformer。MM-DiT 为图像和文字保留各自的参数变换，在联合注意力中交换信息，使文字表征也能随图像状态更新。这与固定文字表征供图像交叉查询的接口不同。\cite{a_sd3}

用简化的矩阵写法，先分别计算
\[
(Q_I,K_I,V_I)=\mathcal P_I(H_I),\qquad
(Q_T,K_T,V_T)=\mathcal P_T(H_T),
\]
再在注意力操作中拼接键、查询和值。例如图像侧更新使用
\[
\operatorname{softmax}\!\left(
 \frac{Q_I[K_I;K_T]^\top}{\sqrt{d_k}}\right)[V_I;V_T].
\]
文字侧也可以读取图像侧信息。分号表示沿序列维度拼接。两种模态不用强迫所有变换共享一套参数，但能够在每层协调。这是“多模态”一词在这里的具体含义：两类 token 共同参与场的计算；它不自动意味着模型能回答任意视觉问题。

在直线路径的\hyperref[sec:flowmatching]{流匹配}设定下，按本书时钟写为
\[
\ell_\tau=(1-\tau)Z+\tau\ell_0,
\qquad u^{\mathrm{target}}=\ell_0-Z.
\]
网络回归这个条件速度的最优平均，再由 ODE 生成隐变量。这里的速度与第\ref{sec:parameterization}节的 $v$-prediction 不是同一个定义，不能只因为两者都在代码里叫 \texttt{v} 就相互替换。原论文可能使用从数据到噪声的时钟，读者还需检查符号方向。

图\ref{fig:application-interfaces}把三种接口放在一起。所有箭头表示信息通道，不代表所有组件都端到端一起训练。能认清这张图，阅读新的生成模型时就不必被名字牵着走：先找它的表示空间，随后找模态之间在哪里交流，再找它究竟预测什么。

\begin{figure}[htbp]
\centering
\begin{tikzpicture}[x=1cm,y=1cm,every node/.style={font=\small},
 b/.style={draw,rounded corners,align=center,minimum height=0.65cm},
 ar/.style={->,thick}]
\node[anchor=east] at (-0.1,0) {交叉注意力};
\node[b,minimum width=2.1cm] (ta) at (1.7,0) {固定文字特征};
\node[b,minimum width=2.5cm] (ia) at (5.6,0) {含噪图像主干};
\draw[ar] (ta)--(ia);
\node[anchor=east] at (-0.1,-1.4) {联合注意力};
\node[b,minimum width=2.1cm] (tb) at (1.7,-1.4) {文字状态};
\node[b,minimum width=2.5cm] (ib) at (5.6,-1.4) {图像状态};
\draw[<->,thick] (tb)--node[above]{每层交流}(ib);
\node[anchor=east] at (-0.1,-2.8) {统一任务};
\node[b,minimum width=2.1cm] (tc) at (1.7,-2.8) {文字与视觉上下文};
\node[b,minimum width=2.5cm] (ic) at (5.6,-2.8) {共享或分专家主干};
\draw[ar] (tc)--(ic);
\node[b,minimum width=2.5cm] (out) at (5.6,-4.2) {文字概率／视觉速度};
\draw[ar] (ic)--(out);
\end{tikzpicture}
\caption{条件接口逐步走向统一任务的三种概念结构。双向交流只是结构能力，仍需任务和数据让它形成有用协作。}
\label{fig:application-interfaces}
\end{figure}

\section{编辑与控制：在可能的画面中选择符合要求的一张}
\label{sec:application-control}

\subsection{为什么一句文字不足以规定一幅图}

“一个人在街上骑自行车”可以对应无数姿态、构图和视角。文字擅长说出语义，却很难高效指定每条轮廓。条件图像、深度图、边缘图和姿态骨架能够补充空间约束。数学上，任务从 $p(x\mid c_{\mathrm{text}})$ 变成
\[
p(x\mid c_{\mathrm{text}},c_{\mathrm{spatial}}).
\]
新增条件使允许的画面集合缩小，但通常仍然保留大量纹理、光照和背景自由度。控制生成的目的，就是在约束下保留合理变化，而不是把所有自由度消灭。

补全是一个直观例子。令掩码 $M$ 在已知区域取一，在需要生成的区域取零，已知像素为 $x_{\mathrm{known}}$。一个理想约束是
\[
M\odot x_0=M\odot x_{\mathrm{known}}.
\]
生成器要根据可见区域推断不可见区域的条件分布。采样过程若持续使用已知区域的相应噪声状态，能够加强一致性；具体方法也可以训练专门接收掩码和条件图的网络。掩码边界附近仍需协调纹理、光照和几何，简单粘贴并不能代替条件建模。

\subsection{ControlNet：给预训练先验增加空间输入}

ControlNet 冻结预训练生成主干，用可训练分支处理附加空间条件，并通过零初始化卷积将分支信号送回主干。它让较小的任务数据能够利用已有图像先验，支持边缘、深度和姿态等条件。\cite{a_controlnet}

零初始化在这里有一个清楚的作用。训练刚开始，附加分支对主干的注入接近零，因此系统行为接近原来的生成器；随着训练进行，条件分支逐渐学会对不同尺度提供修正。这是一种有用的工程策略：不要让一个尚未学会含义的新接口立即扰乱成熟的表征。

条件图和生成图并非像素对像素的硬复制。骨架固定大致关节位置，仍然允许不同衣服和人物；深度固定相对空间结构，仍然允许不同材质；边缘图本身有歧义，同一轮廓可能对应不同对象。控制效果要评价约束是否被满足，也要评价保留的自由度是否合理。

\subsection{DreamBooth 与指令编辑：记住主体和理解修改}

DreamBooth 通过少量主体照片微调文字到图像模型，把一个特别的标识与主体关联，并用类别先验保留思想缓解微调后泛化受损的问题。它研究的是“这个具体对象如何进入已有生成先验”。\cite{a_dreambooth}

主体个性化的难点是区分身份与偶然环境。一组照片若都在同一背景、同一光照下拍摄，模型可能把背景当成身份的一部分。换一个地点生成时，便可能连背景也一起带走。增加视角和环境变化，并适度限制微调幅度，才有机会学到更稳定的主体特征。这个问题本质上是条件下的可辨识性：训练数据是否足以告诉模型哪些变化属于身份，哪些只是拍摄情境。

InstructPix2Pix 则学习原图、编辑指令和目标图之间的关系，利用合成的编辑配对训练条件扩散模型，推理时可直接依据指令修改图像。它强调编辑行为需要专门的监督，而不是仅凭文字生成能力便会自然出现。\cite{a_instructpix}

对编辑任务，单纯评价最终图像漂亮与否是不够的。若要求“把杯子变成蓝色”，模型同时改变桌面、人物和构图，就没有正确完成局部指令。应当至少检查三件事：要求修改的属性是否改变，主体身份是否保留，不相关区域是否稳定。生成质量和编辑忠实度并不总是同向变化。

\begin{cautionbox}{编辑强度与保持程度的矛盾}
从原图加更多噪声再生成，通常会给模型更大改动自由，但也削弱原图约束。改得少不一定更正确，改得多也不一定更听话。理想编辑需要理解“应该改变什么”和“必须保持什么”，而不只是一个统一的强度旋钮。
\end{cautionbox}

\section{视频生成：时间一致性不是逐帧画图}
\label{sec:application-video}

\subsection{从边缘分布到联合分布}

设视频有 $F$ 帧，单独为每帧生成一张好图，只能保证每个边缘分布看起来合理。视频需要的是联合分布
\[
p(x^{(1)},\ldots,x^{(F)}\mid c).
\]
同一个人物跨帧应保持身份，物体运动应连续，遮挡后重现的东西应保留先前属性。即使每一帧单独看都毫无问题，它们的组合仍可能出现衣服闪烁、物体漂移和人数跳变。这与概率论里“边缘分布不能唯一确定联合分布”完全一致，相关概念见附录\ref{app:probability}。

早期 Video Diffusion Models 用时空网络研究视频的联合去噪，并探索图像和视频共同训练。它是理解从二维生成到时序生成的重要入口：时间轴进入网络，学习目标也随之涵盖跨帧相关性。\cite{a_videodiff}

可以把图像 U-Net 中的空间算子扩展成时空算子，也可以分别做空间注意力和时间注意力。后一种分解中，同一帧内的 token 先交换信息，相同空间位置附近的 token 再跨帧交流。它降低直接在全部时空 token 上交互的压力，但分解的顺序和范围也影响模型对复杂运动的表达。

\subsection{时空压缩和视频 Transformer}

视频张量比图像多了时间维度，内存与计算问题更加突出。若直接把 $F$ 帧都切成 token，数量约为 $N=Fhw/p^2$，全注意力中的 $N^2$ 项会随帧数平方增长。时空自编码器因此不仅压缩空间，也压缩时间，把多帧映射成较短的隐变量序列。

Wan 的 2025 年技术报告公开了视频 DiT、时空 VAE、数据处理与训练策略。其 Wan-VAE 使用因果时间结构；报告说明普通后续帧可按时间四倍、空间八倍压缩，第一帧单独保留时间处理方式。这些是具体公开实现，应与产品系列其他版本分开。\cite{a_wan}

“因果”在这里指一个时间位置的计算不使用未来帧信息，它支持分块编码、解码和缓存设计；它不意味着整个视频生成主干只能按帧自回归。自编码器的因果性、Transformer 的注意力范围与采样的时间顺序，属于三种不同层面的选择。

一个重要取舍是时间压缩率。压缩过强时，短促动作和嘴唇运动可能在接口中丢失；压缩较弱时，长视频计算昂贵。把时空压缩只当作无损的加速器，会低估它对运动质量的影响。对于视频也应先做重建实验：真实视频经过编解码后，快速运动和细小文字是否仍然稳定？

\subsection{视频生成与世界模型之间还隔着什么}

OpenAI 2024 年 Sora 技术披露说明了压缩视觉数据、时空 patches 和扩散 Transformer 的路线。披露同时展示了一些持续性与模拟能力，也报告了物理和长时间一致性的局限。它适合作为公开研究思路的案例，无法据此还原全部训练细节。\cite{a_sora2024}

一段看起来真实的物体运动，不等于模型已经掌握可供规划的真实动力学。视频数据告诉模型哪些画面经常相邻，却可能没有告诉它哪些变化由动作造成。要把模型用于决策，还需要回答反事实问题：如果施加另一个动作，后续状态如何变化？需要区分相机移动与物体移动，需要在新的环境中保持预测校准，还需要为不确定性留下空间。

例如，一只玻璃杯在许多视频里被完整拿起。模型据此生成“完整拿起”的画面很合理；但若机器人用过大的力抓杯子，真实结果可能是破裂。仅从常见画面学习到的统计规律，不一定覆盖关键而罕见的动作后果。一个能够辅助规划的世界模型，需要在行动条件下学习动态，而不仅是生成视觉上顺畅的故事。

Sora 2 的 2025 年官方材料披露同步音视频生成能力，但没有给出足以独立重建系统的完整网络和训练目标。本书将它作为产品能力的历史例子，不把其演示推断为特定未公开架构。按官方后续公告，Sora 网页与应用在 2026 年 4 月 26 日停止，API 公告停止日期为 2026 年 9 月 24 日；这个状态变化也说明电子书应区分技术贡献和服务可用性。\cite{a_sora2,a_soraend}

\section{统一多模态：理解与生成怎样共享计算}
\label{sec:application-unified}

\subsection{“能看图”和“能画图”为什么不自动统一}

视觉理解需要从像素中提取适合回答问题的语义。例如，回答“桌上有几个杯子”时，背景纹理大多不重要，物体数量却非常重要。视觉生成需要恢复可解码的细节，包括理解任务愿意忽略的纹理、颜色和局部位置。一种理解表征可能对光照变化保持不变，这对识别有帮助；但生成器如果需要改变光照，便不能把所有光照信息都抹去。

因此，统一多模态不只是把所有东西转成 token，然后宣布问题解决。它需要考虑至少三个接口：用于读取图像的表征，用于生成图像的表征，以及二者与文字上下文之间的交流方式。一个系统可以为理解和生成使用不同视觉编码器，同时在主干中共享上下文；也可以学习共同隐变量，再用不同解码器得到不同模态。

文本的离散概率目标与图像的连续回归目标也不同。一个常见混合训练的概念形式是
\begin{align}
\mathcal L={}&\lambda_T\,
\E[-\log p_\theta(w_j\mid w_{<j},\text{上下文})]\\
&+\lambda_I\,
\E[\norm{u_\theta(\ell_\tau,\tau,\text{上下文})
-(\ell_0-Z)}^2].
\label{eq:application-hybrid-loss}
\end{align}
本式是教学抽象，不是某一模型完整训练配方。$\lambda_T,\lambda_I$ 调节任务权重；$w_j$ 是第 $j$ 个文字 token。两项对共享参数提出不同要求，任务比例、数据质量和参数分工都会影响协作还是干扰。

\subsection{Transfusion：共享主干，保留不同生成机制}

Transfusion 把文字的下一 token 预测与图像的扩散目标放入同一个多模态模型，使用离散文字与连续图像表征，在图像块内部允许双向交互。这说明统一主干不要求所有模态都使用同一种输出分布。\cite{a_transfusion}

为什么图像内部的双向交流合理？图像的整个含噪状态在一次去噪调用时已经可见，我们要联合预测它的更新，没有必要人为规定“左上角必须先于右下角”。文字生成却通常需要遵守因果顺序，只让当前位置读取已经确定的前文。统一模型因而需要按模态调整注意力可见范围。

这是一条容易遗漏的训练原则。若文字目标能看到未来正确答案，训练误差会虚假降低，推理却不能复现；若图像被迫只有单向局部上下文，又会丢失扩散并行协同的一部分优势。注意力掩码必须与生成时实际可见的信息一致。

\subsection{BAGEL：两种视觉表征与专家分工}

BAGEL 的 2025 年研究采用文字自回归与视觉 Rectified Flow 的组合，使用理解用的视觉编码器和生成用的 VAE 表征。其 MoT 结构对理解与生成 token 使用不同参数专家，在统一注意力中交流，并用大规模交错多模态数据训练。\cite{a_bagel}

这里的“专家”不是几个会互相聊天的小模型。它指计算图中不同 token 类型使用不同的可训练参数。以这种方式，理解任务可保留语言与语义处理的能力，生成任务可拥有专门处理连续视觉状态的变换，而上下文仍能跨模态传播。总参数量、每个 token 激活的参数量与一次完整生成的计算量，是不同概念，不能混着比较。

交错数据指一条样本中可能出现文字、图像、再文字、再图像。例如，先给出人物照片，再说明要更换衣服，然后给出修改结果，随后继续要求改变背景。它与孤立的“图片加标题”配对相比，明确教模型如何把之前的视觉内容当成后续任务的上下文。系统想获得多轮编辑、跨图身份一致性和图文交替表达，就需要相应的数据组织，而不只是更多无关图片。

读者可以用一个小例子检验自己的理解：第一轮让模型画一只戴红围巾的猫，第二轮要求“让它坐在窗边，保留围巾”。第二轮条件不应只有一句孤立文字，还应包含第一轮的视觉信息。否则“它”指的是哪只猫，以及围巾的具体样子，都缺少稳定依据。多模态上下文提供记忆接口，模型是否有效利用仍需实际评估。

\subsection{2026 年的两个研究窗口}

MUNI 是 2026 年 6 月的预印本，研究共享随机隐变量下的任意模态到任意模态生成。它把模态编码器、解码器与共享流先验联合训练，并讨论跨模态一致性、子集条件的预测充分性和隐变量简洁性之间的关系。其理论结论依赖可实现性和优化等假设。\cite{a_muni}

这项研究引出的一个基础问题值得深入思考。假设我们要同时生成“雨中的街景图”和“雨声”，分别生成一张好街景、一个好雨声，只保证各自边缘分布。若图中雨很小，声音却是猛烈暴雨，它们没有表达同一个情境。共享隐变量希望承载共同场景，再由模态解码器表达各自细节。可共享信息太少，会出现不一致；共享信息太多，又可能把声音里的随机相位、图片里的随机纹理都强塞进同一个复杂变量。这是一个关于充分性与压缩的概率建模问题。

WorldBagel 是 2026 年 7 月的预印本，把 BAGEL 式统一模型扩展到视觉、语言、动作和未来观察，在机器人任务上研究统一训练的作用。它提出动作表征和交错任务组织，并报告若干指定环境中的实验；这些实验不能当作通用物理模拟已经实现的证据。\cite{a_worldbagel}

这里最有启发的是任务关系的改变。过去，文字到图像问“给定描述，什么画面合理”；行动世界建模问“给定当前观察和动作，什么后果合理”。它必须把动作放进条件，考虑多步误差与决策反馈。生成器面对的对象也从静态图像分布扩展到环境转移分布。两个预印本作为正在发展的研究方向阅读，比把它们当成已定型标准更合适。

\subsection{商业多模态系统：明确公开知识的边界}

OpenAI 在 2025 年发布的 4o 图像生成材料强调多模态上下文、文字渲染和编辑能力；2026 年 9 月 8 日的 ChatGPT Images 2.5 发布材料继续强调细节、精确编辑与多轮保持。这里我们引用的是官方产品披露，不能从表现推导出完整主干网络、噪声路径或训练损失。\cite{a_gptimage,a_images25}

同一个能力可以由不同组合实现：语言模型先形成语义计划再交给图像生成器，共享 Transformer 同时处理文字和连续视觉状态，或者通过其他尚未公开的混合设计。界面里的“原生多模态”说明的是系统层面的能力和使用方式，不是一条足够精确的数学定义。

评价这样的系统，可以研究可观察行为：多轮修改是否保留身份，文本信息是否准确，复杂关系是否满足，延迟是否符合任务需求。解释底层原理则应依赖论文、技术报告、公开代码和可复现实验。把这两类证据分开，既能欣赏系统实际进展，也能保持技术判断清楚。

\section{从声音到分子：扩散学习的不一定是像素}
\label{sec:application-beyond}

\subsection{音频：同一个条件，多种合理声音}

音频生成可在波形、时频谱或压缩音频隐变量中进行。AudioLDM 研究文本到音频的隐空间扩散，用音频表征建立生成条件，并把压缩表征解码为声音。它说明“压缩表示加条件生成”的思想可以从图像移植到别的感知模态。\cite{a_audioldm}

一条提示“雨点落在金属屋顶上”，不会规定每个雨点的准确落下时刻。因此，输出仍然是条件分布中的一个样本。音频的挑战包括持续时间、事件先后、节奏和声源一致性。若任务是语音，音素内容还必须精确；若任务是音乐，长时间结构与和声关系很重要。只让短窗口内的声音真实，并不能保证整段作品的组织合理。

时频谱把声音分解为时间和频率附近的能量信息，但从谱回到波形还需处理相位或依赖专门解码器。因而“把音频当成一张图片来生成”只是一种局部直觉，不能忽略声音特有的物理和时间结构。

\subsection{三维与分子：先选对几何空间}

三维对象可以表示为点云、网格、体素、隐式场或可渲染场景。表示一换，网络要遵守的结构也换了。点云对点的排列通常不敏感，三维旋转不应改变对象本身；网格还要考虑连接关系。第\ref{ch:math}章中直接向所有欧氏坐标加高斯噪声的例子，是一个重要起点，但不一定是所有几何任务的合适终点。

DiffDock 把分子对接写成配体位姿的生成问题，在平移、旋转和扭转自由度构成的空间上建立扩散过程，并配合置信度估计。关键变化是：生成变量不是一幅分子图片，而是实际几何自由度。\cite{a_diffdock}

旋转矩阵受到正交性和行列式约束，不能把矩阵的每个元素任意扰动后，仍当作合法旋转。周期性的扭转角也不能简单按无边界实数处理。这样的例子提醒我们，score、噪声和速度必须尊重变量所在空间的几何；术语可参考\hyperref[gl:manifold]{流形}与\hyperref[gl:equivariance]{等变性}条目。

RFdiffusion 将蛋白质结构设计表述为逐步去噪的生成任务，并研究基序支架、结合体等条件设计。论文包含实验验证的设计案例，因此展示了从生成结构到物理实验验证的完整研究链条。\cite{a_rfdiffusion}

这里的质量标准明显不同于图像。一个结构外观像蛋白质，并不意味着它会折叠成预期形态，更不意味着它具有所需功能。计算预测、结构筛选和实验测量分别提供不同证据。对于科学应用，生成器提出候选，验证流程决定哪些候选可信。扩散模型扩展了可以探索的设计空间，却没有替代领域中的约束与检验。

\section{怎样评价一种进展}
\label{sec:application-evaluation}

\subsection{一个分数不能承担所有结论}

\term{FID}{gl:fid}在固定特征空间里比较生成与真实样本的统计量，它对分布偏差有用，却不直接回答一张图是否符合某条复杂指令。图文相似度可以检查总体语义，却可能漏掉数量、左右关系和文字拼写。人类偏好评测更贴近视觉体验，但结果会受到提示集合、评价问题、样本选择和比较对手影响。

表\ref{tab:application-evaluation}按任务而非模型列出评价问题。它是一份原创检查框架，不是论文里的性能排行榜。一个实际实验应先选定任务，再确定哪些错误最有代价，然后选择能显露这些错误的指标。

\begin{table}[htbp]
\centering
\small
\begin{tabularx}{\linewidth}{@{}p{2.1cm}XX@{}}
\toprule
任务 & 最该检查的问题 & 单一视觉质量分数容易漏掉什么\\
\midrule
文字到图像 & 对象、数量、属性和关系是否同时满足 & 正确构图中的错误数量或绑定\\
局部编辑 & 指令是否完成，其他部分是否保持 & 漂亮但更换了主体和背景\\
空间控制 & 生成是否遵守骨架、深度或边缘 & 纹理逼真但几何偏离\\
视频 & 身份、运动、遮挡和长期状态是否连续 & 每帧清晰但跨帧闪烁\\
多模态生成 & 各模态是否表达同一事件与时间结构 & 图片与声音各自合理却不匹配\\
科学设计 & 约束、可实现性与实验功能是否成立 & 看似合理的结构不能实际工作\\
\bottomrule
\end{tabularx}
\caption{从任务出发选择评价问题。评价应观察完整条件任务，而不只是输出是否好看。}
\label{tab:application-evaluation}
\end{table}

一个公平的生成比较，还需要报告采样预算。模型一次生成一张图，与先生成二十张再筛选最好的一张，体现的是不同计算预算。采样步数、随机种子数、提示词重写、引导强度和后处理都应清楚说明。论文展示一组很好的图，只能证明系统能产生这些图；稳定成功率需要在有代表性的提示集合上验证。

\subsection{用消融实验定位贡献}

消融实验是有控制地改变一个组件，观察结果的实验。若新系统同时更换文字编码器、扩大数据、增加模型宽度和改变损失，最终变好，不能直接说全部收益来自新损失。一个有说服力的结论需要尽可能固定其他条件，或者明确承认组件间耦合。

读论文时，可沿着一条具体问题链阅读：作者先测量了什么失败？新方法改变哪个组件？对照是否只有该组件不同？收益在哪些任务成立？代价是更多训练还是更多推理？是否公开了失败样例？这些问题比“模型名字里有没有最新术语”更能帮助判断一个进展。

\section{一条可走通的论文阅读路线}
\label{sec:application-reading}

下面选择二十篇具有明确教学价值的论文或技术报告。它们并非二十个竞争榜单位置，而是二十个不同的问题窗口。第二、三章的 DDPM、DDIM、score-SDE 和 Flow Matching 奠定数学前提；读到本章时，优先看问题设置、方法图和消融表，再回到推导。每条入口直接链接一手来源。

\begin{enumerate}[leftmargin=*,itemsep=0.65em]
\item \textbf{ADM，2021。}问题：扩散的图像先验怎样与强生成基线竞争？关注结构消融和引导的质量、覆盖率取舍；局限是类别条件并不等同于开放语言控制。\href{https://arxiv.org/abs/2105.05233}{阅读入口}。\cite{a_adm}
\item \textbf{GLIDE，2021。}问题：怎样在去噪过程中遵循文字？关注条件表示、引导比较和补全设置；不要把公开的小模型表现当成全部论文实验。\href{https://arxiv.org/abs/2112.10741}{阅读入口}。\cite{a_glide}
\item \textbf{Imagen，2022。}问题：图像模型之外，还有什么限制图文一致性？优先读文本编码器与扩容实验，再看级联生成；复杂关系仍需要专门评估。\href{https://arxiv.org/abs/2205.11487}{阅读入口}。\cite{a_imagen}
\item \textbf{DALL-E 2，2022。}问题：语义表征怎样成为中间生成变量？先画清先验和解码器的分工；语义接口会压缩部分细节。\href{https://arxiv.org/abs/2204.06125}{阅读入口}。\cite{a_dalle2}
\item \textbf{LDM，2021／2022。}问题：高分辨率扩散能否在更紧凑空间中进行？先检查自编码器重建，再读条件机制；压缩损失构成能力边界。\href{https://arxiv.org/abs/2112.10752}{阅读入口}。\cite{c_rombach2022ldm}
\item \textbf{DreamBooth，2022。}问题：怎样让已有先验记住少量照片中的主体？关注身份与类别泛化；数据少时容易把背景学成身份。\href{https://arxiv.org/abs/2208.12242}{阅读入口}。\cite{a_dreambooth}
\item \textbf{InstructPix2Pix，2022。}问题：怎样学习“照着指令改图”？关注训练配对从何而来以及保持能力；合成配对的偏差会影响可学编辑。\href{https://arxiv.org/abs/2211.09800}{阅读入口}。\cite{a_instructpix}
\item \textbf{DiT，2022／2023。}问题：Transformer 能否成为可扩展的扩散主干？重点读条件注入和计算规模实验；原始类别实验不能直接推出所有文图任务。\href{https://arxiv.org/abs/2212.09748}{阅读入口}。\cite{a_dit}
\item \textbf{ControlNet，2023。}问题：怎样利用成熟先验满足空间条件？追踪可训练分支怎样进入主干；约束服从与视觉质量需要分别评估。\href{https://arxiv.org/abs/2302.05543}{阅读入口}。\cite{a_controlnet}
\item \textbf{AudioLDM，2023。}问题：隐空间条件生成怎样迁移到声音？关注条件表征与声音解码接口；长时间事件组织不能由短片段质量替代。\href{https://arxiv.org/abs/2301.12503}{阅读入口}。\cite{a_audioldm}
\item \textbf{SDXL，2023。}问题：如何把隐空间系统推进到更高分辨率？看尺寸条件和组件分工；多项改动需要消融定位收益。\href{https://arxiv.org/abs/2307.01952}{阅读入口}。\cite{a_sdxl}
\item \textbf{SD3／MM-DiT，2024。}问题：如何让文字和图像在流模型里共同计算？区分路径实验与结构实验；比较损失大小时注意目标和权重不同。\href{https://arxiv.org/abs/2403.03206}{阅读入口}。\cite{a_sd3}
\item \textbf{Video Diffusion Models，2022。}问题：从图像边缘分布走向视频联合分布，需要改变什么？关注时空交互与联合训练；短时一致性不等于长期动力学。\href{https://arxiv.org/abs/2204.03458}{阅读入口}。\cite{a_videodiff}
\item \textbf{Wan，2025。}问题：开放的视频基础模型怎样组织编解码、数据和训练？报告较长，可先读模型设计；具体版本与评测条件必须一起记录。\href{https://arxiv.org/abs/2503.20314}{阅读入口}。\cite{a_wan}
\item \textbf{Transfusion，2024。}问题：自回归文字与连续视觉生成怎样共享主干？先看注意力掩码和两个目标；统一参数仍面对任务平衡。\href{https://arxiv.org/abs/2408.11039}{阅读入口}。\cite{a_transfusion}
\item \textbf{BAGEL，2025。}问题：大规模交错数据怎样支持理解与生成？关注专家、两种视觉表征和上下文；任务能力应与数据组织一同讨论。\href{https://arxiv.org/abs/2505.14683}{阅读入口}。\cite{a_bagel}
\item \textbf{MUNI，2026，预印本。}问题：多模态共同生成怎样保持一致？先读潜变量标准，再检查理论假设；基准结果不能替代所有模态上的证明。\href{https://arxiv.org/abs/2606.16408}{阅读入口}。\cite{a_muni}
\item \textbf{WorldBagel，2026，预印本。}问题：统一模型怎样加入动作和未来状态？关注转移任务与动作表征；指定机器人环境中的收益有清楚适用范围。\href{https://arxiv.org/abs/2607.03461}{阅读入口}。\cite{a_worldbagel}
\item \textbf{DiffDock，2022／2023。}问题：存在多个合理配体位姿时，怎样生成并筛选？关注几何空间和置信度；结构误差与实际药物功能属于不同证据。\href{https://arxiv.org/abs/2210.01776}{阅读入口}。\cite{a_diffdock}
\item \textbf{RFdiffusion，2023。}问题：怎样在条件下生成可验证的蛋白质结构？先看设计任务及实验链条；生成候选仍需结构与功能检验。\href{https://www.nature.com/articles/s41586-023-06415-8}{阅读入口}。\cite{a_rfdiffusion}
\end{enumerate}

读完其中任意一篇，可以尝试用五句话复述：它生成什么随机变量，如何构造训练样本，网络如何接收条件，推理如何更新状态，实验排除了哪一种替代解释。若五句话还说不清，继续查结构图和任务设置，通常比再看十张示例图更有帮助。

\section{未来方向：从已知瓶颈提出判断}
\label{sec:application-future}

以下是基于本章讨论的研究判断，不是关于下一代产品的已知事实，也不是保证会实现的时间表。

第一，\textbf{更准确的表示接口}可能与更大的主干同样重要。小字、快速动作、几何细节若在编码器中丢失，后面的场学习无法稳定补救。未来的表示可能按内容分配不同粒度，使平坦背景用较少 token，文字和关键动作获得更细的表征。但自适应接口还必须保持可训练、可解码和高效批处理，不能只在概念上节省 token。

第二，\textbf{生成计算可能更依赖任务难度}。简单背景也许不需要与复杂空间关系同样多的推理预算。第三章的蒸馏、少步生成和求解器改进提供了加速基础；另一条路线是让模型检查约束，发现失败后继续修改。难点在于检查器也会犯错，而且多次生成再选择的额外成本应当诚实计算。

第三，\textbf{多轮编辑需要可维护的视觉状态}。一幅图连续修改十次，哪些对象身份必须保持，哪些属性已经被用户改变，不能仅靠一个不断增长的文字摘要猜测。场景结构、对象引用和连续视觉表征之间的配合，可能成为更可靠的编辑接口。但显式结构并非越多越好：过硬的接口也可能妨碍风格化和创作自由。

第四，\textbf{多模态统一将继续面对共享与专门化的取舍}。所有模态共用一套参数，结构简洁却可能产生任务干扰；为每个模态完全独立建模，又浪费共同知识。共享注意力、专门参数、不同输出分布和统一上下文的组合，是一个开放设计空间。真正有意义的进展应当证明新增协作改善了任务，而不是仅仅让结构图看起来统一。

第五，\textbf{世界建模需要行动和可检验的因果约束}。视频生成带来丰富视觉先验，但用于机器人和科学预测时，误差不再只是画面瑕疵。未来系统需要识别分布外状态，表达多种后果，在闭环中接受真实反馈。把动态模型当成一个允许检验和修正的预测工具，会比把漂亮视频直接等同于真实世界理解更有成效。

这些方向最终又回到本书的第一句话：生成是在学习分布。应用越丰富，分布的对象、条件和评价越具体。扩散与流模型提供了一个强大框架，真正完成任务还需要正确的表示、可学习的条件、合适的数据，以及诚实的验证。

\section{练习与解答思路}
\label{sec:application-exercises}

\begin{exercisebox}{练习 1：定位小字错误}
某文字到图像系统在 $1024\times1024$ 图像中总把商品标签的小字画成乱码。设计三个实验，分别判断错误是否来自自编码器、文字条件或生成先验。
\end{exercisebox}

\begin{exercisebox}{练习 2：计算 token 的代价}
隐变量尺寸为 $64\times64\times16$。将 patch 边长从 $4$ 改成 $2$，token 数、注意力矩阵元素数各改变多少倍？如果图像空间边长再翻倍且压缩率不变，这两项又改变多少倍？说明为什么这些倍率不等于完整推理时间倍率。
\end{exercisebox}

\begin{exercisebox}{练习 3：联合分布与逐帧质量}
构造一个最小例子：每帧都有同样的合理边缘分布，但两种视频生成器的时间一致性完全不同。解释为什么只计算单帧指标不能区分它们。
\end{exercisebox}

\begin{exercisebox}{练习 4：为统一模型设计掩码}
一条训练样本依次包含文字说明、输入图像、编辑指令和目标图像。目标图像用于速度回归，后续文字用于自回归回答。说明哪些 token 应看见哪些信息，并指出一种会导致训练时信息泄漏的设计。
\end{exercisebox}

\begin{exercisebox}{练习 5：把未来方向写成实验}
选择“多轮身份保持”“图像和声音一致”或“动作条件世界预测”之一，提出一个有控制的实验。给出训练条件、比较基线、评价指标和能够推翻你假设的结果。
\end{exercisebox}

\paragraph{练习 1 的思路。}
先对真实标签图做编解码重建，检查字形是否保留；再用内容只有一个大字、一个短词和一条明确关系的提示检验条件理解；最后固定编码器与文字编码器，比较数据或生成主干变体。若重建就失败，先验实验无法单独解决接口缺失。评价文字可用转录准确率与人工核验，不只用整体图像相似度。

\paragraph{练习 2 的思路。}
token 数从 $256$ 变成 $1024$，增加四倍；注意力矩阵元素数增加十六倍。固定 patch 后空间边长翻倍，token 数又增加四倍，注意力矩阵又增加十六倍。线性投影、前馈层、自编码器、采样调用次数和硬件利用率也影响时间，因此完整延迟需要测量。

\paragraph{练习 3 的思路。}
令每帧只有红球或蓝球，概率各为 $1/2$。系统 A 每帧独立选色；系统 B 在第一帧选色后一直保持。所有单帧边缘分布相同，但 A 频繁闪烁，B 始终一致。可用相邻帧颜色改变率区分，说明联合指标不可由单帧指标替代。

\paragraph{练习 4 的思路。}
目标含噪图像可以读取先前说明、输入图像和编辑指令，并在目标图像内部双向交流。文字预测应遵守其实际推理顺序。若让当前文字读取尚未生成的正确未来文字，便发生直接泄漏；若后续任务在训练中读取不可能在推理时获得的干净目标表示，也会形成条件分布不一致。需先明确推理顺序，再设计掩码。

\paragraph{练习 5 的思路。}
以多轮编辑为例，固定主干、训练预算和提示集合，比较有交错数据与仅孤立配对数据的系统；记录每轮指令成功率、主体保持和无关区域变化。若加入交错数据后只让画面更好看，却没有提高保持或任务完成率，原假设未获支持。报告全部轮次和失败，而不只选最佳结果。

\chaptertakeaway{应用发展的主线不是名字不断变新，而是接口不断扩展：U-Net 协调不同图像尺度，隐空间降低处理维度，DiT 把视觉场写成 token 交互，空间条件使生成可控，视频要求联合时间建模，统一多模态让理解与生成共享上下文。无论生成像素、声音、位姿还是蛋白质，都应先选对随机变量与几何空间，再决定条件、网络和验证方法。}
